%!TEX root = ../main.tex

\begin{abstract}
\noindent
\textbf{Background:}
Ubiquitous computing represents the cornerstone of the Ecological Momentary Interventions (EMIs) and Just-in-Time Adaptive Interventions (JITAIs) for mental health. However, this paradigm emphasizes internal physiological sensing (symptom tracking) over the external contextual factors (situational causes), even though it's directly correlated with the mental issue.\newline

\noindent
\textbf{Objectives:}
The objective of this study is to systematically review, categorize, and analyze how physical, social, and computing/temporal contexts are integrated into consumer wearable and smartphone-based mental health interventions.\newline

\noindent
\textbf{Eligibility Criteria:}
Included peer-reviewed articles from journals and conferences evaluated mental health interventions utilizing consumer-grade wearables or smartphones with a primary focus on adaptive, external context sensing. This study, clinical-grade diagnostics, or exclusively tracking internal physiological metrics without external environmental correlation, were excluded.\newline

\noindent
\textbf{Sources:}
A comprehensive search was conducted across Scopus, PsycINFO, PubMed, and Web of Science.\newline

\noindent
\textbf{Selection Methods:}
Following PRISMA guidelines, titles and abstracts were screened using automated machine learning tools(ASReviwer), followed by independent full-text reviews. Data was extracted utilizing a standardized Physical, Social, and Computing framework.\newline

\noindent
\textbf{Results:}
112 studies met the final inclusion criteria. The synthesis reveals that [...]. \newline

\noindent
\textbf{Conclusions:}
The shift from passive symptom monitoring to active, situation-aware intervention models shows an enhanced ecological validity of digital mental health tools. Bridging the existing "why" gap in ubiquitous computing, future research should consider prioritizing multimodal external sensing while targeting a robust, privacy-preserving computational frameworks.\newline

\end{abstract}